\documentclass{article}
\usepackage{amsmath, amssymb}
\usepackage{graphicx} % Required for inserting images

\title{DL- Homework 1}
\author{Mrinal Bharadwaj}
\date{January 2026}

\begin{document}

\maketitle


\section*{Question 1 (Optimization)}

Compute the gradient $\nabla f(x,y)$ and Hessian $\nabla^2 f(x,y)$ of the function
\[
f(x,y)=-(x^2-1)^2-y^2+2xy.
\]
Find all stationary points of this function. Show that
\[
\left(\sqrt{\frac{3}{2}},\sqrt{\frac{3}{2}}\right)
\]
is a local maximum.

\section*{Answer to Question 1}

First, expand the function:
\[
f(x,y)=-x^4+2x^2-1-y^2+2xy.
\]

\subsection*{Gradient}
\[
\frac{\partial f}{\partial x}=-4x^3+4x+2y,
\qquad
\frac{\partial f}{\partial y}=2x-2y.
\]

\[
\nabla f(x,y)=
\begin{bmatrix}
-4x^3+4x+2y\\
2x-2y
\end{bmatrix}.
\]

\subsection*{Hessian}
\[
\frac{\partial^2 f}{\partial x^2}=4-12x^2,
\quad
\frac{\partial^2 f}{\partial y^2}=-2,
\quad
\frac{\partial^2 f}{\partial x\partial y}=2.
\]

\[
\nabla^2 f(x,y)=
\begin{bmatrix}
4-12x^2 & 2\\
2 & -2
\end{bmatrix}.
\]

\subsection*{Stationary Points}

Solve $\nabla f(x,y)=0$.

From $2x-2y=0$, we get $y=x$.
Substitute into the first equation:
\[
-4x^3+6x=0
\quad\Rightarrow\quad
2x(2x^2-3)=0.
\]

Thus,
\[
x=0 \quad \text{or} \quad x=\pm\sqrt{\frac{3}{2}}.
\]
Since $y=x$, the stationary points are:
\[
(0,0),\;
\left(\sqrt{\frac{3}{2}},\sqrt{\frac{3}{2}}\right),\;
\left(-\sqrt{\frac{3}{2}},-\sqrt{\frac{3}{2}}\right).
\]

\subsection*{Local Maximum}

At $x^2=\frac{3}{2}$, the Hessian becomes:
\[
\nabla^2 f=
\begin{bmatrix}
-14 & 2\\
2 & -2
\end{bmatrix}.
\]

Compute principal minors:
\[
D_1=-14<0,
\qquad
D_2=\det(\nabla^2 f)=28-4=24>0.
\]

Hence, the Hessian is negative definite, and
\[
\left(\sqrt{\frac{3}{2}},\sqrt{\frac{3}{2}}\right)
\]
is a strict local maximum.

\newpage

\section*{Question 2 (Optimization)}

Show that the function
\[
f(x,y)=x^3-3xy^2
\]
has only one stationary point, and that it is neither a local minimum nor a local maximum, but a saddle point.

\section*{Answer to Question 2}

\subsection*{Stationary Point}

Compute the gradient:
\[
\frac{\partial f}{\partial x}=3x^2-3y^2,
\qquad
\frac{\partial f}{\partial y}=-6xy.
\]

Set $\nabla f(x,y)=0$:
\[
3(x^2-y^2)=0 \Rightarrow x^2=y^2,
\]
\[
-6xy=0 \Rightarrow xy=0.
\]

The only solution satisfying both equations is:
\[
(x,y)=(0,0).
\]

\subsection*{Classification}

Evaluate $f$ along different paths approaching $(0,0)$.

Along $y=0$:
\[
f(x,0)=x^3,
\]
which is positive for $x>0$ and negative for $x<0$.

Along $y=x$:
\[
f(x,x)=x^3-3x^3=-2x^3,
\]
which is negative for $x>0$ and positive for $x<0$.

Thus, in every neighborhood of $(0,0)$, the function takes both positive and negative values relative to $f(0,0)=0$.

\subsection*{Conclusion}
\[
(0,0)\text{ is neither a local minimum nor a local maximum, but a saddle point.}
\]

\newpage
\section*{Question 3 (Linear Algebra)}

If $A$ and $B$ are positive definite matrices, prove that the block diagonal matrix
\[
\begin{bmatrix}
A & 0 \\
0 & B
\end{bmatrix}
\]
is also positive definite.

\section*{Answer to Question 3}

Recall that a symmetric matrix $M$ is positive definite if and only if
\[
x^\top M x > 0 \quad \text{for all } x \neq 0.
\]

Let
\[
M =
\begin{bmatrix}
A & 0 \\
0 & B
\end{bmatrix},
\qquad
z =
\begin{bmatrix}
x \\
y
\end{bmatrix},
\]
where $x \in \mathbb{R}^n$ and $y \in \mathbb{R}^m$, not both zero.

Then
\[
z^\top M z
=
\begin{bmatrix}
x^\top & y^\top
\end{bmatrix}
\begin{bmatrix}
A & 0 \\
0 & B
\end{bmatrix}
\begin{bmatrix}
x \\
y
\end{bmatrix}
=
x^\top A x + y^\top B y.
\]

Since $A$ and $B$ are positive definite,
\[
x^\top A x > 0 \quad \text{for } x \neq 0,
\qquad
y^\top B y > 0 \quad \text{for } y \neq 0.
\]

If $z \neq 0$, then at least one of $x$ or $y$ is nonzero, implying
\[
z^\top M z = x^\top A x + y^\top B y > 0.
\]

Therefore, it
\[
\begin{bmatrix}
A & 0 \\
0 & B
\end{bmatrix}
\]
is positive definite.
\newpage
\section*{Question 4 (Chain Rule Calculus)}

Consider the function
\[
f(x)=w_2^\top \,\sigma(W_1 x),
\]
where $\sigma(\cdot)=\text{sigmoid}(\cdot)$ applies element-wise to a vector and
\[
\sigma(t)=\frac{1}{1+e^{-t}}, \qquad \sigma'(t)=\sigma(t)\bigl(1-\sigma(t)\bigr).
\]
Assume $W_1\in\mathbb{R}^{c\times d}$, $x\in\mathbb{R}^{d\times 1}$, and $w_2\in\mathbb{R}^{c\times 1}$.
Compute the derivatives $\frac{\partial f}{\partial w_2}$, $\frac{\partial f}{\partial W_1}$, and $\frac{\partial f}{\partial x}$.

\section*{Answer to Question 4}

Define the intermediate variables
\[
z = W_1 x \in \mathbb{R}^{c\times 1}, \qquad s=\sigma(z)\in\mathbb{R}^{c\times 1}.
\]
Then
\[
f = w_2^\top s \in \mathbb{R}.
\]

\subsection*{1) Derivative w.r.t. $w_2$}
Since $f = w_2^\top s$ is linear in $w_2$,
\[
\frac{\partial f}{\partial w_2} = s = \sigma(W_1 x)
\qquad (\in \mathbb{R}^{c\times 1}).
\]

\subsection*{2) Derivative w.r.t. $x$}
First compute $\frac{\partial f}{\partial z}$:
\[
\frac{\partial f}{\partial s} = w_2,
\qquad
\frac{\partial s}{\partial z} = \sigma'(z)= s\odot(1-s),
\]
so (element-wise chain rule)
\[
\frac{\partial f}{\partial z}
=
\frac{\partial f}{\partial s}\odot \frac{\partial s}{\partial z}
=
w_2 \odot s \odot (1-s)
\qquad (\in \mathbb{R}^{c\times 1}),
\]
where $\odot$ denotes element-wise (Hadamard) product.

Because $z=W_1 x$, we have
\[
\frac{\partial f}{\partial x}
=
W_1^\top \frac{\partial f}{\partial z}
=
W_1^\top \bigl(w_2 \odot s \odot (1-s)\bigr)
\qquad (\in \mathbb{R}^{d\times 1}).
\]

\subsection*{3) Derivative w.r.t. $W_1$}
Using $z=W_1x$ and the fact that $\frac{\partial z}{\partial W_1}$ produces an outer product,
\[
\frac{\partial f}{\partial W_1}
=
\left(\frac{\partial f}{\partial z}\right) x^\top
=
\bigl(w_2 \odot s \odot (1-s)\bigr) x^\top
\qquad (\in \mathbb{R}^{c\times d}).
\]

\subsection*{Final Results (compact form)}
Let $s=\sigma(W_1x)$. Then
\[
\boxed{\frac{\partial f}{\partial w_2}=s}
\]
\[
\boxed{\frac{\partial f}{\partial x}=W_1^\top\bigl(w_2\odot s\odot(1-s)\bigr)}
\]
\[
\boxed{\frac{\partial f}{\partial W_1}=\bigl(w_2\odot s\odot(1-s)\bigr)\,x^\top.}
\]

\newpage
\section*{Question 5 (High Dimensional Statistics: Curse of Dimensionality)}

Consider $N$ data points independently and uniformly distributed in a
$p$-dimensional unit ball
\[
B=\{x\in\mathbb{R}^p : \|x\|_2 \le 1\}
\]
centered at the origin. The median distance from the origin to the closest
data point is given by
\[
d(p,N)=\left(1-2^{-1/N}\right)^{1/p}.
\]

Prove this expression. Then compute the median distance $d(p,N)$ for
$N=10{,}000$ and $p=1{,}000$.

\section*{Answer to Question 5}

Let $R_{\min}$ denote the distance from the origin to the closest of the
$N$ data points.

\subsection*{Step 1: Probability distribution of $R_{\min}$}

For a point drawn uniformly from the unit ball, the probability that its
distance from the origin is at most $r \in [0,1]$ is proportional to the
$p$-dimensional volume contained within radius $r$. Since the total volume
is normalized to $1$, this probability is
\[
\mathbb{P}(\|x\| \le r) = r^p.
\]

Hence,
\[
\mathbb{P}(\|x\| > r) = 1 - r^p.
\]

Because the $N$ points are independent,
\[
\mathbb{P}(R_{\min} > r) = (1 - r^p)^N.
\]

\subsection*{Step 2: Definition of the median}

By definition, the median distance $d(p,N)$ satisfies
\[
\mathbb{P}(R_{\min} > d(p,N)) = \frac{1}{2}.
\]

Substituting the expression above,
\[
(1 - d(p,N)^p)^N = \frac{1}{2}.
\]

Taking the $N$th root,
\[
1 - d(p,N)^p = 2^{-1/N}.
\]

Solving for $d(p,N)$ gives
\[
d(p,N)^p = 1 - 2^{-1/N},
\]
\[
d(p,N) = \left(1 - 2^{-1/N}\right)^{1/p}.
\]

\subsection*{Step 3: Numerical evaluation}

For $N=10{,}000$ and $p=1{,}000$,
\[
d(1000,10000)
=
\left(1 - 2^{-1/10000}\right)^{1/1000}.
\]

Using the approximation $2^{-1/N} \approx 1 - \frac{\ln 2}{N}$ for large $N$,
\[
d(1000,10000)
\approx
\exp\!\left(\frac{1}{1000}
\ln\!\left(\frac{\ln 2}{10000}\right)\right)
\approx 0.99.
\]

\subsection*{Conclusion}

Even with a large number of data points, in high dimensions the closest
point to the origin lies very close to the boundary of the unit ball,
demonstrating the \textbf{curse of dimensionality}.


\end{document}
